\chapter{Discussion}
\label{ch:discussion}
% BUDGET : 4 pages -- HUIT axes, 0,5 page chacun (cf. memoire/PLAN.md, section 5).
% THESE DU CHAPITRE : ce qui rend ce travail utilisable n'est pas ce qu'il AFFIRME, c'est ce
% qu'il DELIMITE -- et cette delimitation a un cout de position que je peux nommer.
%
% REGLE DE REDACTION, mecanique : chaque section s'ouvre sur le FAIT MESURE, et ensuite
% seulement sur ce qu'il engage. JAMAIS l'inverse. Le risque de ce chapitre n'est pas de ne
% pas etre compris, c'est l'auto-flagellation -- qui est une autre forme de complaisance
% (REF-60). Le fait technique reste le fait technique.

\minisommaire{%
  \ms{sec:disc-stage}{conditions de travail, positionnement, valorisation de la recherche}%
  \ms{sec:disc-ia}{ce que l'assistant de programmation a permis et déplacé}}


Ce chapitre explore certaines réflexions autour du stage et du projet. C'est l'occasion de
discuter de questions de positionnement, de critiquer et de proposer des pistes d'améliorations.

\section{Discussion sur le stage}
\label{sec:disc-stage}

Ce stage a été majoritairement réalisé en autonomie. Je suis arrivé dans la structure début
mai 2026 et mes encadrants sont partis en vacances fin juillet pour un mois. Ils ne sont pas
formés aux outils que je mobilise dans mon stage et je suis assez seul pour le c\oe ur du
travail. J'ai eu une grande liberté en ce qui concerne les choix stratégiques de conception et
mes encadrants ne sont intervenus que ponctuellement.

Nous avons décidé de concentrer le stage sur la partie programmation. L'idée était de dégrossir
cette partie technique en profitant de mes compétences sur le sujet et de réserver la partie
terrain, moins niche, à de futurs stages ou thèses. Le but étant d'avoir un modèle opérationnel
rapidement avec une difficulté liée à la prise en main bien moindre que celle liée à la
programmation. Les prochaines personnes à travailler sur le sujet travailleront à l'exploitation
du modèle, à la campagne de données auprès des agriculteurs ainsi qu'au travail dans la
micro-ferme Karusmart de l'\gls{inrae}. Ce choix se défend également car c'est compliqué de
reprendre en cours un projet de programmation qui n'est pas le nôtre.

Ce stage était donc surtout une affaire d'ordinateur et de lignes de code. C'est quelque chose
auquel je suis habitué et que je commence à maîtriser. Seulement, c'est un cadre de travail qui
n'offre que trop peu de contact humain à mon goût. Je souhaite un travail qui m'amène à
interagir avec d'autres personnes qu'une équipe fixe. Je souhaite de la nouveauté dans mon
quotidien et j'aimerais être moins statique et sédentaire.

Je remarque également que l'approche de modélisation peut avoir quelque chose de froid et
déconnecté de la réalité d'un territoire. Les héritages coloniaux deviennent de simples lignes
de code. C'est le cas pour l'interdiction de certaines cultures sur un sol chlordéconé. Pour
certains, c'est un traumatisme avec des conséquences sanitaires et culturelles. Pour moi, c'est
une contrainte du modèle.

Le problème est dans ce que ce rapport à l'information fait à l'engagement. Un agriculteur
guadeloupéen aurait peut-être une approche beaucoup plus empathique de la situation avec
des conséquences très positives sur les motivations et la détermination de mener le projet.
Cela pourrait peut-être se traduire par une envie à tout prix de créer de l'impact avec ce
projet avec des idées intelligentes qui sortent du cadre académique de recherche classique en
termes de communication, et d'inclusion des acteurs.

Je viens de l'Hexagone, je suis blanc et je n'ai pas de lien avec le monde de l'agriculture ou
de l'aide alimentaire. On pourrait rétorquer que j'ai les compétences pour le stage et que ça
suffit mais le fait que je sois de passage en Guadeloupe, sans lien durable avec le territoire
crée un désengagement qui n'aurait peut-être jamais existé si je travaillais sur ma région
natale et sur un sujet qui me touche personnellement.

Pour remédier à ça, recruter des gens localement qui sont touchés par la question semble idéal.
À défaut de trouver des profils adaptés, je pense qu'il faudrait plus lier les chercheurs aux
acteurs du territoire afin de créer du lien et de la confiance. Cela permettrait d'ajouter une
dimension sensible au travail qu'on effectue et ça aurait pour conséquence de penser
différemment la place du chercheur.

Cette étape de concertation et d'ateliers participatifs est censée arriver après mon stage. Ce
sera l'occasion pour les acteurs de donner leur avis sur le travail réalisé, de la démarche, de
témoigner de leurs difficultés et de participer à la suite du projet.

Par ailleurs, le projet \gls{tifig} ne réalise pas, faute de moyen, d'évaluation d'impact.
Malgré une volonté parmi les membres du laboratoire de faire de la science transformative, peu
de moyens sont mis dans la création de débouchés. La recherche reste descriptive car elle n'est
pas adossée à une méthode d'initiation et de catalysation des changements (rapport de force
médiatico-politique, inclusion forte des acteurs du territoire et des financeurs du projet,
etc.). En bref, il manque un pôle de valorisation de la recherche et les moyens qui vont avec.

Pour conclure, c'était un stage qui m'a plu techniquement mais qui manque selon moi
d'interactions humaines et qui délaisse trop l'approche sensible de l'engagement au profit d'une
approche analytique. Je pense que, par son fonctionnement, la recherche amène vers ce cadre de
pensée et de travail. Je pense vouloir désormais m'orienter vers le travail associatif et
explorer d'autres manières de s'engager professionnellement, rencontrer des personnes avec un
profil différent de ceux que l'on trouve habituellement en recherche et, avec ça, cultiver un
autre rapport au monde.

\section{Discussion sur l'IA générative}
\label{sec:disc-ia}

J'ai utilisé Claude Code tout au long de ce stage. C'est un \gls{llm} spécialisé dans le
développement de code. L'outil a permis un portage de \num{\mesLignesCode} lignes,
\num{\mesTests} tests, \num{\mesChecksums} sommes de contrôle et un dispositif de prospection
opérationnel en \num{4} mois avec une seule personne. Cet outil change radicalement le rapport à
la programmation. J'avais plus un rôle de concepteur/relecteur lors de mon stage. L'essentiel de
l'écriture est réalisé par l'agent.

J'ai utilisé Claude Code pour créer des tableaux et faire de l'analyse de données. L'agent a
également contribué à la rédaction de mon brouillon de travail et m'a permis de mettre
correctement en forme mon rapport sous format \LaTeX.

C'est une aide très intéressante pour accélérer les actions répétitives comme la recherche
de fautes d'orthographe ou le formatage du code. C'est également utile pour la réalisation de
tâches complexes mais peu intéressantes comme l'implémentation d'un \emph{dashboard} d'affichage
des résultats. Là où il faut être vigilant en tant qu'utilisateur, c'est de ne pas déléguer tout
à l'agent. Il présente plein de défauts d'analyse et peut produire des résultats faux mais
plausibles et faire des sauts logiques. De plus, une densification trop rapide du code fait
perdre la maîtrise du projet à celui qui travaille dessus. Le programme peut vite devenir une
boîte noire et Claude Code devient vite indispensable pour toute modification future. Cela crée
une dépendance d'autant plus forte que les compétences de l'utilisateur n'évoluent pas aussi
rapidement que le code se complexifie.

En conclusion, Claude Code est un outil formidable pour toutes sortes de tâches informatiques,
mais la perte de compétence ramenée à la complexité de la tâche, et la délégation excessive de
choix conceptuels et de l'implémentation à Claude Code crée une relation de dépendance à l'agent
à laquelle il faut être vigilant.

À noter que le processus de rédaction de ce mémoire m'a obligé à plonger dans le code et m'a
permis de mieux aligner mon niveau de compréhension du sujet avec le niveau de complexité du
programme et des données produites. C'est bien de faire régulièrement le bilan du travail
en cours pour ne pas perdre le contrôle du projet.
